\documentclass[11pt]{article}
\usepackage{mathblatt}
% Prüfstein 08.10.2026 abends, als Lernblatt (Lehrer 09.10.): „Hypotenuse berechnen“ (Pythagoras, Lerneinheit 1), von oben gebaut.
% Lernweg (vom Lehrer bestätigt 08.10.): 1 Seiten benennen · 2 Satz in Worten · 3 Gleichung aufstellen
% ohne Rechnen · 4 Rechnen glatt, a) grau vorgerechnet, Probe · 5 Rechenaufwand kurz · 6 Sache mit
% wechselnder Frage · 7 zwei P10-Zielaufgaben. Präambel aus durchgaenge-pythagoras.
% Gemeinsame Präambel der drei Durchgänge, übernommen aus paket-pythagoras-2/2-blatt.tex (Probe 08.10.2026).
\makeatletter
\setlength{\mbpfbe}{0mm}% keine Punkte (Bauregel 6.4)
% eingerückt (Bauregel 4.7, 6.6): \gEin Einzug, \gSchrift eine Stufe kleiner
\newlength{\gEin}\setlength{\gEin}{0pt}
\let\gSchrift\relax
\renewcommand{\pfkreuzzeile}[1]{\par\noindent\hangindent3.2em\hangafter1\hspace*{1.6em}$\square$\ #1\par}
\newcommand{\gkreuz}[1]{$\square$\ #1\hspace{2em}}
% Bauregel 6.7: links, was man ansieht, rechts, was man tut; Spaltengrenze überall an derselben Stelle
\newsavebox{\gbox}\newlength{\gLinks}
\newcommand{\gzwei}[2]{\par\smallskip\noindent\sbox{\gbox}{#1}%
  \setlength{\gLinks}{\dimexpr0.5\textwidth-\gEin-\mbpfnr\relax}%
  \ifdim\wd\gbox>\dimexpr\gLinks-4mm\relax
    \usebox{\gbox}\par\smallskip\noindent\begin{minipage}{\linewidth}#2\end{minipage}%
  \else
    \begin{minipage}[t]{\gLinks}\vspace{0pt}\usebox{\gbox}\end{minipage}%
    \begin{minipage}[t]{\dimexpr\linewidth-\gLinks\relax}\vspace{0pt}\raggedright #2\end{minipage}%
  \fi\par}
\RenewDocumentEnvironment{pfaufg}{m m m}{%
  \begin{lrbox}{\mbaufgabebox}\begin{minipage}[t]{\linewidth}%
  \gSchrift\noindent\hspace*{\gEin}\mbpfrandmarke{#1}%
  \begin{minipage}[t]{\mbpfnr}\strut\textbf{#2}\end{minipage}%
  \begin{minipage}[t]{\dimexpr\linewidth-\gEin-\mbpfnr\relax}\raggedright\strut\ignorespaces}%
 {\end{minipage}%
  \end{minipage}\end{lrbox}%
  \setlength{\mbaufgabehoehe}{\dimexpr\ht\mbaufgabebox+\dp\mbaufgabebox\relax}%
  \par\addvspace{7pt}\Needspace*{\mbaufgabehoehe}%
  \noindent\usebox{\mbaufgabebox}\par\addvspace{7pt}}
\makeatother
% Rechenraum als Karo (Bauregel 6.9): n Rechenschritte = n mal 10 mm
\newcommand{\karo}[1]{\par\smallskip\noindent\begin{tikzpicture}
  \clip (0,0) rectangle ({\linewidth-1mm},{#1*10mm});
  \draw[black!16,line width=0.3pt,step=5mm] (0,0) grid ({\linewidth},{#1*10mm});
  \end{tikzpicture}\par}
\newcommand{\kk}[1]{$\square$\,#1\hspace{0.9em}}
\newcommand{\linie}[1]{\,{\color{black!45}\rule[-1pt]{#1}{0.4pt}}\,}
% Zeile am Ende jedes Durchgangs: Originale klein grau, darüber; dann Vorgänger/Nachfolger (Bauregel 3.4, 6.5)
\newcommand{\blattende}[2]{\par\vfill\noindent\begin{minipage}{\linewidth}{\scriptsize\color{mbgrau}Originale: #1\par}\smallskip
{\footnotesize #2\par}\end{minipage}}
\newcommand{\pkt}{\textperiodcentered{}}

\newcommand{\kennung}{K4W}
\newcommand{\grau}[1]{{\color{mbgrau}#1}}

\begin{document}
\pfheftstil
\fancyfoot[C]{\parbox[t]{\textwidth}{\mbfhbox\par\vspace{1.5mm}{\footnotesize\color{mbgrau}{\scriptsize \kennung}\hfill\thepage}}}
\pfheftkopf{Hypotenuse berechnen}{Klasse 9}

% 1 Seiten benennen: erst Standardlagen, dann gedreht (Fehler: längere Kathete als Hypotenuse)
\begin{pfaufg}{}{1.}{}
Gegenüber dem rechten Winkel liegt die längste Seite des Dreiecks. Sie heißt \textbf{Hypotenuse}. Die beiden Seiten am rechten Winkel heißen \textbf{Katheten}.\par
Kreuze in jedem Dreieck die Hypotenuse an.\par\medskip\noindent
\begin{minipage}[b]{0.245\linewidth}\centering
\begin{tikzpicture}[line width=0.7pt,scale=0.85]\coordinate (P1) at (0,0);\coordinate (P2) at (2.4,0);\coordinate (P3) at (0,1.5);\draw (P1) -- (P2) -- (P3) -- cycle;\mbrwmarke{(P1)}{(P2)}{(P3)}
\node[font=\small,anchor=north] at (1.2,-0.08) {$a$};\node[font=\small,anchor=east] at (-0.08,0.75) {$b$};\node[font=\small,anchor=south west] at (1.22,0.78) {$c$};\end{tikzpicture}\par\smallskip
{\small a)\ \kk{$a$}\kk{$b$}\kk{$c$}}
\end{minipage}\hfill
\begin{minipage}[b]{0.245\linewidth}\centering
\begin{tikzpicture}[line width=0.7pt,scale=0.85]\coordinate (P1) at (0,0);\coordinate (P2) at (2.4,0);\coordinate (P3) at (2.4,1.5);\draw (P1) -- (P2) -- (P3) -- cycle;\mbrwmarke{(P2)}{(P1)}{(P3)}
\node[font=\small,anchor=north] at (1.2,-0.08) {$p$};\node[font=\small,anchor=west] at (2.48,0.75) {$q$};\node[font=\small,anchor=south east] at (1.18,0.78) {$r$};\end{tikzpicture}\par\smallskip
{\small b)\ \kk{$p$}\kk{$q$}\kk{$r$}}
\end{minipage}\hfill
\begin{minipage}[b]{0.245\linewidth}\centering
\begin{tikzpicture}[line width=0.7pt,scale=0.85]\coordinate (P1) at (0,0);\coordinate (P2) at (2.5,0);\coordinate (P3) at (0.9,1.2);\draw (P1) -- (P2) -- (P3) -- cycle;\mbrwmarke{(P3)}{(P1)}{(P2)}
\node[font=\small,anchor=north] at (1.25,-0.08) {$s$};\node[font=\small,anchor=south east] at (0.40,0.62) {$t$};\node[font=\small,anchor=south west] at (1.75,0.62) {$u$};\end{tikzpicture}\par\smallskip
{\small c)\ \kk{$s$}\kk{$t$}\kk{$u$}}
\end{minipage}\hfill
\begin{minipage}[b]{0.245\linewidth}\centering
\begin{tikzpicture}[line width=0.7pt,scale=0.85]\coordinate (P1) at (0,1.25);\coordinate (P2) at (1.8,0);\coordinate (P3) at (2.8,1.44);\draw (P1) -- (P2) -- (P3) -- cycle;\mbrwmarke{(P2)}{(P1)}{(P3)}
\node[font=\small,anchor=north east] at (0.85,0.56) {$f$};\node[font=\small,anchor=north west] at (2.37,0.67) {$g$};\node[font=\small,anchor=south] at (1.40,1.42) {$e$};\end{tikzpicture}\par\smallskip
{\small d)\ \kk{$e$}\kk{$f$}\kk{$g$}}
\end{minipage}
\fusshilfe{\mbox{1:~a) $c$; b) $r$; c) $s$; d) $e$}}
\end{pfaufg}

% 2 Grund des Satzes, nach dem Benennen: Flächen addieren, nicht Längen (Fehler c = a + b; Lehrer 09.10.)
\begin{pfaufg}{}{2.}{}
Über jeder Seite des rechtwinkligen Dreiecks steht ein Quadrat. Die Hypotenuse ist 5 Kästchen lang.
\gzwei{\begin{tikzpicture}[scale=0.40,line width=0.7pt]
\draw[mbgitter,line width=0.3pt] (-4,-3) grid (7,7);
\draw (0,0) -- (3,0) -- (3,-3) -- (0,-3) -- cycle;
\draw (0,0) -- (-4,0) -- (-4,4) -- (0,4) -- cycle;
\draw (0,4) -- (3,0) -- (7,3) -- (4,7) -- cycle;
\draw[line width=1.4pt] (0,0) -- (3,0) -- (0,4) -- cycle;
\mbrwmarke{(0,0)}{(3,0)}{(0,4)}
\end{tikzpicture}}{%
\textbf{a)}\ Zähle die Kästchen in jedem Quadrat.\par\smallskip
\hspace*{1.4em}Katheten:\linie{12mm}und\linie{12mm}\par\smallskip\hspace*{1.4em}Hypotenuse:\linie{12mm}\par\bigskip
\textbf{b)}\ Kreuze an, was stimmt.
\pfkreuzzeile{Die zwei Katheten sind zusammen so lang wie die Hypotenuse.}
\pfkreuzzeile{Die zwei kleinen Quadrate haben zusammen so viele Kästchen wie das große.}}
\fusshilfe{\mbox{2:~a) 9, 16, 25; b) die Quadrate}}
\end{pfaufg}

% 3 Gleichung aufstellen ohne Rechnen: wechselnde Buchstaben, Punktnamen wie in der P10,
% ein Dreieck ohne rechten Winkel (Fehler: „c ist immer die Hypotenuse“, Satz ohne rechten Winkel)
\begin{pfaufg}{}{3.}{}
Das gilt in jedem rechtwinkligen Dreieck. Es ist der Satz des Pythagoras: \mbox{Hypotenuse$^2$ = Kathete$^2$ + Kathete$^2$.}\par\smallskip
Schreibe für jedes Dreieck die Gleichung des Satzes auf. Hat ein Dreieck keinen rechten Winkel, schreibe „gilt nicht“.\par
Suche zuerst die Hypotenuse – sie steht allein auf einer Seite der Gleichung.\par\medskip\noindent
\begin{minipage}[b]{0.245\linewidth}\centering
\begin{tikzpicture}[line width=0.7pt,scale=0.85]\coordinate (P1) at (0,0);\coordinate (P2) at (2.5,0);\coordinate (P3) at (1.6,1.2);\draw (P1) -- (P2) -- (P3) -- cycle;\mbrwmarke{(P3)}{(P1)}{(P2)}
\node[font=\small,anchor=north] at (1.25,-0.08) {$w$};\node[font=\small,anchor=south east] at (0.78,0.62) {$u$};\node[font=\small,anchor=south west] at (2.08,0.62) {$v$};\end{tikzpicture}\par\smallskip
{\small a)}\linie{28mm}
\end{minipage}\hfill
\begin{minipage}[b]{0.245\linewidth}\centering
\begin{tikzpicture}[line width=0.7pt,scale=0.85]\coordinate (P1) at (0,0);\coordinate (P2) at (2.4,0);\coordinate (P3) at (2.4,1.5);\draw (P1) -- (P2) -- (P3) -- cycle;\mbrwmarke{(P2)}{(P1)}{(P3)}
\node[font=\small,anchor=north] at (1.2,-0.08) {$z$};\node[font=\small,anchor=west] at (2.48,0.75) {$y$};\node[font=\small,anchor=south east] at (1.18,0.78) {$x$};\end{tikzpicture}\par\smallskip
{\small b)}\linie{28mm}
\end{minipage}\hfill
\begin{minipage}[b]{0.245\linewidth}\centering
\begin{tikzpicture}[line width=0.7pt,scale=0.85]\coordinate (A) at (0,0);\coordinate (B) at (2.4,0);\coordinate (C) at (0,1.5);\draw (A) -- (B) -- (C) -- cycle;\mbrwmarke{(A)}{(B)}{(C)}
\node[font=\small,anchor=north east] at (A) {$A$};\node[font=\small,anchor=north west] at (B) {$B$};\node[font=\small,anchor=south east] at (C) {$C$};\end{tikzpicture}\par\smallskip
{\small c)}\linie{28mm}
\end{minipage}\hfill
\begin{minipage}[b]{0.245\linewidth}\centering
\begin{tikzpicture}[line width=0.7pt,scale=0.85]\coordinate (P1) at (0,0);\coordinate (P2) at (2.5,0);\coordinate (P3) at (0.35,1.5);\draw (P1) -- (P2) -- (P3) -- cycle;
\node[font=\small,anchor=north] at (1.25,-0.08) {$m$};\node[font=\small,anchor=east] at (0.10,0.75) {$k$};\node[font=\small,anchor=south west] at (1.48,0.78) {$l$};\end{tikzpicture}\par\smallskip
{\small d)}\linie{28mm}
\end{minipage}
\fusshilfe{\mbox{3:~a) $w^2 = u^2 + v^2$; b) $x^2 = y^2 + z^2$; c) $\overline{BC}^2 = \overline{AB}^2 + \overline{AC}^2$; d) gilt nicht}}
\end{pfaufg}

% 4 Rechnen mit glatten Zahlen; a) grau vorgerechnet (Bauregel 1.5), b) selbst; Probe statt Kästchenbild
% (Fehler: Wurzel vergessen, c = a + b)
\begin{pfaufg}{}{4.}{}
Die Katheten eines rechtwinkligen Dreiecks sind gegeben. Berechne die Hypotenuse $c$.\par\medskip
\grau{a)\ $6$ cm und $8$ cm:\quad $c^2 = 6^2 + 8^2 = 36 + 64 = 100$\quad $\Rightarrow$\quad $c = \sqrt{100} = 10$ cm}\par\medskip
\noindent\begin{minipage}[t]{0.48\linewidth}
b)\ $5$ cm und $12$ cm
\karo{3}
\end{minipage}\hfill
\begin{minipage}[t]{0.48\linewidth}
c)\ $9$ cm und $12$ cm
\karo{3}
\end{minipage}
\par Prüfe dein Ergebnis: Die Hypotenuse ist länger als jede Kathete, aber kürzer als beide zusammen.
\fusshilfe{\mbox{4:~b) 13 cm; c) 15 cm}}
\end{pfaufg}

% 5 Rechenaufwand erst nach allen Denkschritten, kurz; Komma am Ende
\begin{pfaufg}{}{5.}{}
Berechne die Hypotenuse $c$ aus den beiden Katheten. Runde, wo nötig, auf eine Stelle nach dem Komma.\par\smallskip
\noindent\begin{minipage}[t]{0.48\linewidth}
a)\ $4$ cm und $9$ cm
\karo{3}
\end{minipage}\hfill
\begin{minipage}[t]{0.48\linewidth}
b)\ $1{,}2$ m und $50$ cm
\karo{3}
\end{minipage}
\fusshilfe{\mbox{5:~a) 9,8 cm; b) 130 cm}}
\end{pfaufg}

% 6–8 Sache: der Schüler findet das Dreieck selbst – erst im Bild ohne Dreieck (Leiter, rechter Winkel
% nur zwischen Wand und Boden), dann auf einer Karte (Strecke selbst einzeichnen), dann nur Text (Skizze
% selbst); Fragen wechseln: reicht es? · wie viel kürzer? · wie lang? (Lehrer 09.10.: Skizzen sahen gleich aus)
\begin{pfaufg}{}{6.}{}
Herr Kaya will das Fenster putzen. Seine Leiter ist $2{,}6$ m lang. Reicht sie bis zum Fenster?
\gzwei{\begin{tikzpicture}[line width=0.7pt,scale=1.05]
\fill[black!10] (3,0) rectangle (3.6,3.6);\draw (3,0) -- (3,3.6);
\draw (2.85,2.4) rectangle (3,3.3);\draw (2.85,2.4) -- (3.6,2.4);
\draw (-0.4,0) -- (4.0,0);
\draw[line width=1.8pt] (2.3,0) -- (2.97,2.4);
\draw[|-|,line width=0.4pt] (2.3,-0.35) -- (3,-0.35) node[midway,below,font=\small] {$0{,}7$ m};
\draw[|-|,line width=0.4pt] (3.9,0) -- (3.9,2.4) node[midway,right,font=\small] {$2{,}4$ m};
\node[font=\small,anchor=east] at (2.45,1.3) {Leiter};
\end{tikzpicture}}{%
\karo{3}\par\smallskip\hfill\kk{ja}\kk{nein}}
\fusshilfe{\mbox{6:~ja, gebraucht werden 2,5 m}}
\end{pfaufg}

\begin{pfaufg}{}{7.}{}
Lena geht von $A$ über die Ecke $K$ nach $C$. Quer durch den Park ist der Weg kürzer. Wie viele Meter spart sie?
\gzwei{\begin{tikzpicture}[line width=0.7pt,scale=0.8]
\fill[black!10] (0,0.45) rectangle (3.55,3.2);
\draw (-0.3,0) -- (4.0,0);\draw (-0.3,0.45) -- (3.55,0.45) -- (3.55,3.2);\draw (4.0,0) -- (4.0,3.2);
\node[font=\scriptsize] at (1.7,1.9) {Park};
\fill (0,0.22) circle (1.6pt) node[font=\small,anchor=east] {$A$};
\fill (3.78,0.22) circle (1.6pt) node[font=\small,anchor=north west] {$K$};
\fill (3.78,3.0) circle (1.6pt) node[font=\small,anchor=south] {$C$};
\draw[|-|,line width=0.4pt] (0,-0.35) -- (3.78,-0.35) node[midway,below,font=\small] {$80$ m};
\draw[|-|,line width=0.4pt] (4.35,0.22) -- (4.35,3.0) node[midway,right,font=\small] {$60$ m};
\end{tikzpicture}}{%
\karo{3}}
\fusshilfe{\mbox{7:~40 m}}
\end{pfaufg}

\begin{pfaufg}{}{8.}{}
Der Bildschirm eines Fernsehers ist $80$ cm breit und $45$ cm hoch. Im Laden wird seine Diagonale angegeben. Berechne sie. Runde auf eine Stelle nach dem Komma.\par
Skizziere zuerst den Bildschirm und zeichne die Diagonale ein.
\karo{4}
\fusshilfe{\mbox{8:~91,8 cm}}
\end{pfaufg}

\clearpage
% 9 Zielaufgabe 1 (2020 · 7 · a, Punkt Q weggelassen): Punktnamen, Zwischenergebnis, Runden
\begin{pfaufg}{nach P10 ’20}{9.}{}
Im Dreieck $ABC$ ist bei $C$ ein rechter Winkel. Die Strecke $\overline{BC}$ ist $12$~m lang, die Strecke $\overline{AC}$ $34$~m. Berechne die Länge der Strecke $\overline{AB}$. Runde auf eine Stelle nach dem Komma.
\gzwei{\begin{tikzpicture}[line width=0.7pt,scale=0.8]\coordinate (A) at (0,0);\coordinate (B) at (1.2,3.4);\coordinate (C) at (0,3.4);\draw (A) -- (B) -- (C) -- cycle;\mbrwmarke{(C)}{(A)}{(B)}\node[font=\small] at (-0.05,-0.3) {$A$};\node[font=\small] at (1.37,3.64) {$B$};\node[font=\small] at (-0.1,3.68) {$C$};\end{tikzpicture}}{%
\karo{3}}
\fusshilfe{\mbox{9:~36,1 m}}
\end{pfaufg}

% 10 Zielaufgabe 2 (2025 · 4 · a, Winkel weggelassen): die kurze Höhe ist keine Hypotenuse
\begin{pfaufg}{nach P10 ’25}{10.}{}
Familie Yücel baut vor der Haustür eine Rampe für Rollstühle. Berechne die Länge $x$ der Rampe. Runde auf eine Stelle nach dem Komma.
\gzwei{\begin{tikzpicture}[line width=0.7pt,scale=0.75]\coordinate (P1) at (0,0);\coordinate (P2) at (5.95,0);\coordinate (P3) at (5.95,0.56);\draw (P1) -- (P2) -- (P3) -- cycle;\mbrwmarke{(P2)}{(P1)}{(P3)}
\node[font=\small,anchor=north] at (2.97,-0.08) {$170$ cm};\node[font=\small,anchor=west] at (6.03,0.28) {$16$ cm};\node[font=\small,anchor=south] at (2.9,0.32) {$x$};\end{tikzpicture}}{%
\karo{3}}
\fusshilfe{\mbox{10:~170,8 cm}}
\end{pfaufg}

\blattende{2020 \pkt{} 7 \pkt{} a \quad 2025 \pkt{} 4 \pkt{} a}{Vorher: Quadrat und Wurzel\quad\textbullet\quad Weiter: Kathete berechnen}
\end{document}
